\documentclass[tikz,border=8pt]{standalone}
\usepackage{ctex}
\usepackage{amsmath}
\usetikzlibrary{calc, arrows.meta, decorations.pathreplacing}

\begin{document}
% 二面角 α-l-β 示意图
% 棱 l + 两半平面 α, β + 平面角标示
\begin{tikzpicture}[
  x={(1cm,0cm)},
  y={(0.42cm,0.30cm)},
  z={(0cm,0.88cm)},
  thick,
]

% 二面角：棱 l 沿 y 轴方向
% 半平面 α 向左上展开，半平面 β 向右下展开

% 棱 l 的两端点
\coordinate (L0) at (0,0,0);
\coordinate (L1) at (0,4,0);

% 平面角取点（在棱上取一点 O'）
\coordinate (Oo) at (0,2,0);

% 半平面 α（左侧，向上）
\coordinate (A1) at (-2.5,0,1.8);
\coordinate (A2) at (-2.5,4,1.8);

% 半平面 β（右侧，向下展）
\coordinate (B1) at (2.5,0,-0.3);
\coordinate (B2) at (2.5,4,-0.3);

% 画半平面 α（蓝色）
\fill[blue!15, draw=blue!50, fill opacity=0.75]
  (L0) -- (A1) -- (A2) -- (L1) -- cycle;
\node[blue!70, font=\small\bfseries] at (-1.8, 3.5, 1.0) {$\alpha$};

% 画半平面 β（绿色）
\fill[green!15, draw=green!60!black, fill opacity=0.75]
  (L0) -- (B1) -- (B2) -- (L1) -- cycle;
\node[green!60!black, font=\small\bfseries] at (2.0, 3.5, -0.1) {$\beta$};

% 画棱 l（红色加粗）
\draw[red, very thick] (L0) -- (L1);
\node[red, font=\small, right] at ($(L1)+(0.05,0.1)$) {$l$};
\fill[red] (L0) circle [radius=1.5pt];
\fill[red] (L1) circle [radius=1.5pt];

% 画平面角：在 O' 处，两半平面各引垂直于 l 的射线 OA', OB'
% α 侧垂线 OA'（在半平面α内，垂直于 l，即沿 x 负方向+z正方向）
\coordinate (OA) at (-1.6, 2, 1.15);  % 在 α 平面内，垂直棱方向
\coordinate (OB) at (1.6, 2, -0.19);  % 在 β 平面内，垂直棱方向

\draw[blue, thick, -Stealth] (Oo) -- (OA);
\node[blue, font=\small] at ($(OA)+(-0.2,0.1)$) {$OA'$};

\draw[green!70!black, thick, -Stealth] (Oo) -- (OB);
\node[green!70!black, font=\small] at ($(OB)+(0.2,0)$) {$OB'$};

% 标记点 O'
\fill[black] (Oo) circle [radius=1.5pt];
\node[font=\small, below right] at ($(Oo)+(0.1,-0.05)$) {$O'$};

% 二面角弧线（在 O' 处，从 OA' 到 OB'）
% 画一段小弧表示角
\draw[orange, thick]
  ($(Oo)!0.45!(OA)$) arc [start angle=145, end angle=345, radius=0.55cm];
\node[orange, font=\small\bfseries] at (0.6, 2, 0.5) {$\theta$};

% 底部说明
\node[font=\small, align=center] at (0, -0.5, 0)
  {二面角 $\alpha$-$l$-$\beta$};
\node[font=\small, align=center] at (0, -0.95, 0)
  {平面角 $\theta = \angle A'O'B' \in [0,\,\pi]$};

\end{tikzpicture}
\end{document}
